\documentclass[stu,letterpaper,12pt,donotrepeattitle,floatsintext]{apa7}
% Preámbulo común de los documentos de Nexus: Normas APA 7.a edición (trabajo de estudiante).
% Se carga con % Preámbulo común de los documentos de Nexus: Normas APA 7.a edición (trabajo de estudiante).
% Se carga con % Preámbulo común de los documentos de Nexus: Normas APA 7.a edición (trabajo de estudiante).
% Se carga con \input{../comun/preambulo} justo después de \documentclass.
\usepackage[spanish,es-noshorthands,es-tabla]{babel}
\usepackage[T1]{fontenc}
\usepackage{newtxtext,newtxmath}
\usepackage{csquotes}
\usepackage[style=apa,backend=biber]{biblatex}
\addbibresource{../comun/referencias.bib}
\usepackage{graphicx}
\usepackage{booktabs}
\usepackage{tabularx}
\usepackage{array}
\usepackage{enumitem}
\usepackage{xurl}
\usepackage{tikz}
\usetikzlibrary{positioning,arrows.meta,shapes.misc}
\hypersetup{hidelinks}
\setlist{itemsep=2pt,topsep=4pt,parsep=0pt}
\setlength{\emergencystretch}{3em}

% Datos de la portada
\authorsnames{Neyder Sebastián Orozco Villamil, Deisy Lorena Jiménez Martínez}
\authorsaffiliations{Universidad Central}
\course{Modelado 3D y Video Juegos (43391312)}
\professor{Mg. Carlos Iván Pinzón Romero}
\duedate{18 de noviembre de 2026}

% Atajos de escritura
\newcommand{\cod}[1]{\texttt{#1}}            % archivos, scripts y código
\newcommand{\ruta}[1]{\textit{#1}}           % rutas de menú: Archivo > Guardar
\newcommand{\tecla}[1]{\textsc{#1}}          % teclas

% Columnas de tabla con texto justificado a la izquierda
\newcolumntype{L}{>{\raggedright\arraybackslash}X}
\newcolumntype{P}[1]{>{\raggedright\arraybackslash}p{#1}}

% Ubicación de figuras y tablas: cerca del texto (los manuales usan H: exactamente en su lugar)
\usepackage{float}
\floatplacement{figure}{!htbp}
\floatplacement{table}{!htbp}
% Nota de figuras y tablas en español y alineada a la izquierda (APA)
\AtBeginDocument{\renewcommand{\notelabel}{Nota}}
\newcommand{\NotaFig}[1]{\if\relax\detokenize{#1}\relax\else\par\raggedright\figurenote{#1}\fi}

% Figura APA: número y título arriba, imagen y nota opcional abajo
%   \Figura[ancho]{archivo}{etiqueta}{título}{nota}
\newcommand{\Figura}[5][0.82\linewidth]{%
  \begin{figure}
    \caption{#4}\label{#3}
    \centering\includegraphics[width=#1]{#2}
    \NotaFig{#5}
  \end{figure}}

% Dos imágenes lado a lado en una misma figura
%   \FiguraDoble{archivo a}{archivo b}{etiqueta}{título}{nota}
\newcommand{\FiguraDoble}[5]{%
  \begin{figure}
    \caption{#4}\label{#3}
    \centering
    \includegraphics[width=0.49\linewidth]{#1}\hfill\includegraphics[width=0.49\linewidth]{#2}
    \NotaFig{#5}
  \end{figure}}

% Tabla APA: número y título arriba; contenido a espacio sencillo y en letra un poco menor
\newenvironment{TablaAPA}[2]{%
  \begin{table}\caption{#2}\label{#1}\centering\small\linespread{1}\selectfont}
  {\end{table}}

% Términos de las referencias como en el manual de normas APA adjunto (Menéndez Echavarría, 2020)
\DefineBibliographyStrings{spanish}{version = {Versión}, nodate = {s.\,f\adddot}}
 justo después de \documentclass.
\usepackage[spanish,es-noshorthands,es-tabla]{babel}
\usepackage[T1]{fontenc}
\usepackage{newtxtext,newtxmath}
\usepackage{csquotes}
\usepackage[style=apa,backend=biber]{biblatex}
\addbibresource{../comun/referencias.bib}
\usepackage{graphicx}
\usepackage{booktabs}
\usepackage{tabularx}
\usepackage{array}
\usepackage{enumitem}
\usepackage{xurl}
\usepackage{tikz}
\usetikzlibrary{positioning,arrows.meta,shapes.misc}
\hypersetup{hidelinks}
\setlist{itemsep=2pt,topsep=4pt,parsep=0pt}
\setlength{\emergencystretch}{3em}

% Datos de la portada
\authorsnames{Neyder Sebastián Orozco Villamil, Deisy Lorena Jiménez Martínez}
\authorsaffiliations{Universidad Central}
\course{Modelado 3D y Video Juegos (43391312)}
\professor{Mg. Carlos Iván Pinzón Romero}
\duedate{18 de noviembre de 2026}

% Atajos de escritura
\newcommand{\cod}[1]{\texttt{#1}}            % archivos, scripts y código
\newcommand{\ruta}[1]{\textit{#1}}           % rutas de menú: Archivo > Guardar
\newcommand{\tecla}[1]{\textsc{#1}}          % teclas

% Columnas de tabla con texto justificado a la izquierda
\newcolumntype{L}{>{\raggedright\arraybackslash}X}
\newcolumntype{P}[1]{>{\raggedright\arraybackslash}p{#1}}

% Ubicación de figuras y tablas: cerca del texto (los manuales usan H: exactamente en su lugar)
\usepackage{float}
\floatplacement{figure}{!htbp}
\floatplacement{table}{!htbp}
% Nota de figuras y tablas en español y alineada a la izquierda (APA)
\AtBeginDocument{\renewcommand{\notelabel}{Nota}}
\newcommand{\NotaFig}[1]{\if\relax\detokenize{#1}\relax\else\par\raggedright\figurenote{#1}\fi}

% Figura APA: número y título arriba, imagen y nota opcional abajo
%   \Figura[ancho]{archivo}{etiqueta}{título}{nota}
\newcommand{\Figura}[5][0.82\linewidth]{%
  \begin{figure}
    \caption{#4}\label{#3}
    \centering\includegraphics[width=#1]{#2}
    \NotaFig{#5}
  \end{figure}}

% Dos imágenes lado a lado en una misma figura
%   \FiguraDoble{archivo a}{archivo b}{etiqueta}{título}{nota}
\newcommand{\FiguraDoble}[5]{%
  \begin{figure}
    \caption{#4}\label{#3}
    \centering
    \includegraphics[width=0.49\linewidth]{#1}\hfill\includegraphics[width=0.49\linewidth]{#2}
    \NotaFig{#5}
  \end{figure}}

% Tabla APA: número y título arriba; contenido a espacio sencillo y en letra un poco menor
\newenvironment{TablaAPA}[2]{%
  \begin{table}\caption{#2}\label{#1}\centering\small\linespread{1}\selectfont}
  {\end{table}}

% Términos de las referencias como en el manual de normas APA adjunto (Menéndez Echavarría, 2020)
\DefineBibliographyStrings{spanish}{version = {Versión}, nodate = {s.\,f\adddot}}
 justo después de \documentclass.
\usepackage[spanish,es-noshorthands,es-tabla]{babel}
\usepackage[T1]{fontenc}
\usepackage{newtxtext,newtxmath}
\usepackage{csquotes}
\usepackage[style=apa,backend=biber]{biblatex}
\addbibresource{../comun/referencias.bib}
\usepackage{graphicx}
\usepackage{booktabs}
\usepackage{tabularx}
\usepackage{array}
\usepackage{enumitem}
\usepackage{xurl}
\usepackage{tikz}
\usetikzlibrary{positioning,arrows.meta,shapes.misc}
\hypersetup{hidelinks}
\setlist{itemsep=2pt,topsep=4pt,parsep=0pt}
\setlength{\emergencystretch}{3em}

% Datos de la portada
\authorsnames{Neyder Sebastián Orozco Villamil, Deisy Lorena Jiménez Martínez}
\authorsaffiliations{Universidad Central}
\course{Modelado 3D y Video Juegos (43391312)}
\professor{Mg. Carlos Iván Pinzón Romero}
\duedate{18 de noviembre de 2026}

% Atajos de escritura
\newcommand{\cod}[1]{\texttt{#1}}            % archivos, scripts y código
\newcommand{\ruta}[1]{\textit{#1}}           % rutas de menú: Archivo > Guardar
\newcommand{\tecla}[1]{\textsc{#1}}          % teclas

% Columnas de tabla con texto justificado a la izquierda
\newcolumntype{L}{>{\raggedright\arraybackslash}X}
\newcolumntype{P}[1]{>{\raggedright\arraybackslash}p{#1}}

% Ubicación de figuras y tablas: cerca del texto (los manuales usan H: exactamente en su lugar)
\usepackage{float}
\floatplacement{figure}{!htbp}
\floatplacement{table}{!htbp}
% Nota de figuras y tablas en español y alineada a la izquierda (APA)
\AtBeginDocument{\renewcommand{\notelabel}{Nota}}
\newcommand{\NotaFig}[1]{\if\relax\detokenize{#1}\relax\else\par\raggedright\figurenote{#1}\fi}

% Figura APA: número y título arriba, imagen y nota opcional abajo
%   \Figura[ancho]{archivo}{etiqueta}{título}{nota}
\newcommand{\Figura}[5][0.82\linewidth]{%
  \begin{figure}
    \caption{#4}\label{#3}
    \centering\includegraphics[width=#1]{#2}
    \NotaFig{#5}
  \end{figure}}

% Dos imágenes lado a lado en una misma figura
%   \FiguraDoble{archivo a}{archivo b}{etiqueta}{título}{nota}
\newcommand{\FiguraDoble}[5]{%
  \begin{figure}
    \caption{#4}\label{#3}
    \centering
    \includegraphics[width=0.49\linewidth]{#1}\hfill\includegraphics[width=0.49\linewidth]{#2}
    \NotaFig{#5}
  \end{figure}}

% Tabla APA: número y título arriba; contenido a espacio sencillo y en letra un poco menor
\newenvironment{TablaAPA}[2]{%
  \begin{table}\caption{#2}\label{#1}\centering\small\linespread{1}\selectfont}
  {\end{table}}

% Términos de las referencias como en el manual de normas APA adjunto (Menéndez Echavarría, 2020)
\DefineBibliographyStrings{spanish}{version = {Versión}, nodate = {s.\,f\adddot}}

\graphicspath{{figuras/}{../desarrollo/figuras/}}
\floatplacement{figure}{H}
\floatplacement{table}{H}

\title{Manual de instalación del videojuego Nexus para Windows}

\begin{document}
\maketitle

\section{Introducción}

\textit{Nexus: Protocolo de emergencia} es un videojuego de terror y sigilo en tercera persona,
desarrollado en Unity \parencite{unity2026} como trabajo final de la asignatura Modelado 3D y Video
Juegos. Este manual explica, paso a paso, cómo instalar y ejecutar la versión para computadores con
Windows, cómo dar los primeros pasos en el juego y cómo resolver los problemas más comunes.

La versión de Windows no necesita un instalador: se distribuye como una carpeta comprimida que
contiene el ejecutable \cod{Nexus.exe} y los archivos que lo acompañan. Para \textquote{instalar} el
juego basta con extraer esa carpeta en cualquier ubicación del equipo.

\section{Requisitos del sistema}

La Tabla~\ref{tab:requisitos} resume los requisitos. El juego se probó en un computador portátil con
Windows~11 y una tarjeta gráfica dedicada; los valores mínimos son una estimación a partir del tamaño
del juego y de la configuración de calidad con que se compiló.

\begin{TablaAPA}{tab:requisitos}{Requisitos del sistema para la versión de Windows}
\begin{tabularx}{\linewidth}{P{3.6cm}LL}
\toprule
Componente & Mínimo & Recomendado \\
\midrule
Sistema operativo & Windows 10 de 64 bits & Windows 11 de 64 bits \\
Procesador & Doble núcleo de 64 bits & Cuatro núcleos o más \\
Memoria RAM & 4 GB & 8 GB \\
Tarjeta gráfica & Compatible con DirectX 11 & Tarjeta dedicada compatible con DirectX 12 \\
Espacio en disco & 250 MB libres & 500 MB libres \\
Controles & Teclado y ratón & Teclado y ratón, o un mando compatible con Windows (XInput) \\
Audio & Altavoces o audífonos & Audífonos, para ubicar a los androides por el sonido \\
\bottomrule
\end{tabularx}
\end{TablaAPA}

\section{Contenido del paquete}

El juego se entrega en el archivo \cod{Nexus\_Windows.zip}, de unos 97 MB. Al extraerlo ocupa unos
217 MB y contiene los elementos de la Tabla~\ref{tab:contenido}. Todos deben permanecer juntos en la
misma carpeta: el ejecutable no funciona si se copia solo.

\begin{TablaAPA}{tab:contenido}{Archivos y carpetas que contiene el paquete de Windows}
\begin{tabularx}{\linewidth}{P{5.2cm}L}
\toprule
Elemento & Función \\
\midrule
\cod{Nexus.exe} & Ejecutable del juego: es el archivo que se abre para jugar \\
\cod{Nexus\_Data} & Escenas, modelos, texturas, sonidos y código del juego \\
\cod{UnityPlayer.dll} & Motor de Unity que ejecuta el juego \\
\cod{MonoBleedingEdge} & Entorno de ejecución del código C\# \\
\cod{D3D12}, \cod{dstorage.dll}, \cod{dstoragecore.dll} & Componentes gráficos y de carga rápida de DirectX \\
\cod{UnityCrashHandler64.exe} & Registra información si el juego se cierra por un error \\
\bottomrule
\end{tabularx}
\end{TablaAPA}

\section{Instalación paso a paso}

\subsection{Paso 1: obtener el archivo comprimido}

Descargue \cod{Nexus\_Windows.zip} o cópielo desde la memoria USB donde se entregó, y guárdelo en una
carpeta del equipo, por ejemplo en \textit{Descargas} o en el \textit{Escritorio}. En el Explorador de
archivos se ve como una carpeta comprimida con un cierre en el ícono (Figura~\ref{fig:zip}).

\Figura[0.85\linewidth]{01_zip_en_carpeta.png}{fig:zip}{El archivo comprimido del juego en el Explorador de archivos}{}

\subsection{Paso 2: extraer los archivos}

Windows puede mostrar el contenido de un archivo \cod{.zip} sin extraerlo, pero el juego no funciona
desde ahí, porque necesita leer sus carpetas como archivos normales. Para extraerlo:

\begin{enumerate}
  \item Haga clic derecho sobre \cod{Nexus\_Windows.zip} y elija \ruta{Extraer todo}.
  \item En el asistente \ruta{Extraer carpetas comprimidas (en zip)} (Figura~\ref{fig:asistente})
        aparece como destino una carpeta con el mismo nombre del archivo. Puede dejarla así o elegir
        otra con \ruta{Examinar}.
  \item Deje marcada la casilla \ruta{Mostrar los archivos extraídos al completar} y pulse
        \ruta{Extraer}. La extracción tarda unos segundos.
\end{enumerate}

\Figura[0.62\linewidth]{02_asistente_extraer.png}{fig:asistente}{Asistente de Windows para extraer la carpeta comprimida}{}

\subsection{Paso 3: abrir la carpeta del juego}

Al terminar, Windows abre la carpeta \cod{Nexus\_Windows} con los archivos de la
Tabla~\ref{tab:contenido} (Figura~\ref{fig:carpeta}). El archivo que se usa para jugar es
\cod{Nexus.exe}, de tipo \textit{Aplicación}.

\Figura[0.85\linewidth]{03_carpeta_extraida.png}{fig:carpeta}{Carpeta del juego después de la extracción}{}

\subsection{Paso 4: ejecutar el juego}

Haga doble clic en \cod{Nexus.exe}. Si el archivo se descargó de Internet, Windows puede mostrar el
aviso \textquote{Windows protegió su PC}. El aviso aparece con las aplicaciones que no tienen una
firma digital comercial y que todavía son poco conocidas para el filtro Microsoft Defender
SmartScreen \parencite{smartscreen}; no significa que el juego tenga un problema. Si aparece:

\begin{enumerate}
  \item Pulse el enlace \ruta{Más información}.
  \item Compruebe que la aplicación es \cod{Nexus.exe} y pulse \ruta{Ejecutar de todas formas}.
\end{enumerate}

Windows recuerda la decisión y el aviso no vuelve a aparecer para ese archivo. En los equipos donde
el juego se copia desde una memoria USB, el aviso normalmente no se muestra.

\subsection{Paso 5: primer inicio}

El juego abre su pantalla de título (Figura~\ref{fig:menu}). Se ejecuta en una ventana que se puede
maximizar o cambiar de tamaño; la combinación \tecla{Alt + Enter} alterna entre ventana y pantalla
completa. Las opciones del menú se eligen con el ratón, con las flechas y \tecla{Enter}, o con el mando.

\Figura[0.85\linewidth]{04_juego_menu.png}{fig:menu}{Pantalla de título del juego ejecutándose en Windows}{}

\subsection{Paso 6 (opcional): crear un acceso directo}

Para abrir el juego desde el escritorio, haga clic derecho sobre \cod{Nexus.exe} y elija
\ruta{Mostrar más opciones} \textrightarrow{} \ruta{Enviar a} \textrightarrow{}
\ruta{Escritorio (crear acceso directo)}. El acceso directo apunta al ejecutable, así que la carpeta
del juego no se debe mover después de crearlo.

\section{Primeros pasos en el juego}

\begin{enumerate}
  \item En el menú elija \ruta{Nueva partida}.
  \item Seleccione a su técnico con las flechas o con los botones \cod{<} y \cod{>}
        (Figura~\ref{fig:selector}): Neyder lanza un señuelo de red que atrae a los androides y Deisy
        apaga un androide durante 10 segundos con el programa KILLSWITCH.exe. Pulse
        \ruta{Confirmar}.
  \item Lea el informe de la misión y pulse \ruta{Iniciar misión}.
  \item El objetivo es recuperar los tres Discos Duros Maestros, uno por nivel, antes de que se
        agote el tiempo de refrigeración y sin que los androides de seguridad lo atrapen.
\end{enumerate}

\Figura[0.85\linewidth]{051_selector-neyder.jpg}{fig:selector}{Selector de personaje}{}

La Tabla~\ref{tab:controles} muestra los controles. El juego usa \textit{controles tipo tanque}:
el personaje avanza hacia donde mira, sin importar el ángulo de la cámara fija.

\begin{TablaAPA}{tab:controles}{Controles con teclado y ratón, y con mando}
\begin{tabularx}{\linewidth}{LP{4.4cm}P{4.2cm}}
\toprule
Acción & Teclado y ratón & Mando \\
\midrule
Avanzar / retroceder & \tecla{W} / \tecla{S} & Stick izquierdo arriba / abajo \\
Girar & \tecla{A} / \tecla{D} & Stick izquierdo a los lados \\
Correr & \tecla{Shift} & Presionar el stick izquierdo \\
Media vuelta rápida & \tecla{Shift + S} & --- \\
Interactuar, recoger y avanzar diálogos & \tecla{E} & Botón Y \\
Habilidad del técnico & \tecla{Q} & Botón X \\
Inventario, archivo y mapa & \tecla{Tab} o \tecla{I} & Start \\
Equipar cuchillo / pistola / escopeta & \tecla{1} / \tecla{2} / \tecla{3} (\tecla{0} guarda el arma) & Cruceta \\
Apuntar & Clic derecho o \tecla{K} & Gatillo izquierdo \\
Disparar o atacar & Clic izquierdo o \tecla{J} & Gatillo derecho \\
Recargar & \tecla{R} & Botón superior derecho \\
Saltar una cinemática & \tecla{Enter}, \tecla{Esc} o \tecla{Espacio} & Start \\
\bottomrule
\end{tabularx}
\end{TablaAPA}

\section{Partidas guardadas}

La partida se guarda en las terminales de respaldo del juego, gastando una memoria USB del
inventario. Al volver a abrir el juego aparece la opción \ruta{Continuar} en el menú. Los datos se
guardan en el registro de Windows del usuario (en la clave
\cod{HKEY\_CURRENT\_USER\textbackslash Software\textbackslash Nexus Team\textbackslash Nexus}), no
dentro de la carpeta del juego; por eso se conservan aunque la carpeta se mueva o se vuelva a
extraer.

\section{Desinstalación}

Como el juego no se instala en el sistema, para quitarlo basta con borrar la carpeta
\cod{Nexus\_Windows} y, si se creó, el acceso directo del escritorio. La partida guardada se puede
borrar también eliminando la clave del registro mencionada en la sección anterior; no es necesario,
porque ocupa muy poco espacio.

\section{Solución de problemas}

La Tabla~\ref{tab:problemas} reúne los problemas más frecuentes y su solución.

\begin{TablaAPA}{tab:problemas}{Problemas frecuentes y cómo resolverlos}
\begin{tabularx}{\linewidth}{P{5cm}L}
\toprule
Problema & Solución \\
\midrule
Aparece \textquote{Windows protegió su PC} & Pulse \ruta{Más información} y luego \ruta{Ejecutar de todas formas} (paso 4) \\
Falta \cod{UnityPlayer.dll} o el juego no abre & El juego se abrió desde dentro del \cod{.zip} o se copió solo el ejecutable. Extraiga la carpeta completa (paso 2) \\
El antivirus pone el juego en cuarentena & Restaure el archivo desde el antivirus y agregue la carpeta del juego como excepción \\
Pantalla negra o cierre al iniciar & Actualice el controlador de la tarjeta gráfica y verifique que el equipo tenga DirectX 11 o superior \\
El juego va lento & Cierre otros programas, juegue en ventana en lugar de pantalla completa y conecte el portátil a la corriente \\
No hay sonido & Revise el dispositivo de salida predeterminado de Windows y el volumen de la aplicación en el mezclador \\
El personaje no responde a las teclas & Haga clic dentro de la ventana del juego para que reciba el teclado \\
\bottomrule
\end{tabularx}
\end{TablaAPA}

\printbibliography

\end{document}
